%% kapa: exam=1 round=37 year=2026 period=2 form=A subject=accounting subject_ko=회계학 seq=3 number=43
%% source: Q-Net 감정평가사 기출문제 (cst00309) artlSeq=5258924 file=raw/1008/5258924/2026년 제37회 감정평가사 1차 2교시.pdf
%% conversion: kordoc 4.24.51 -> tools/gen1.py (auto); review=proofread-sonnet; flags=none
\begin{question}{43}
재고자산에 관한 설명으로 옳지 않은 것은?
\begin{choices}
  \item 후속 생산단계에 투입하기 전에 보관이 필요한 경우 이외의 보관원가는 재고자산의 취득원가에 포함할 수 없으며 발생기간에 비용으로 인식한다.
  \item 순실현가능가치는 재고자산의 주된 또는 가장 유리한 시장에서 시장참여자 사이에 일어날 수 있는 정상거래의 가격에서 처분부대원가를 뺀 금액으로 측정하기 때문에 기업특유의 가치가 아니다.
  \item 생물자산에서 수확한 농림어업 수확물로 구성된 재고자산은 공정가치에서 처분부대원가를 뺀 금액으로 측정하여 수확시점에 최초로 인식한다.
  \item 재고자산을 후불조건으로 취득할 수도 있는데 계약이 실질적으로 금융요소를 포함하고 있다면, 해당 금융요소는 금융이 이루어지는 기간 동안 이자비용으로 인식한다.
  \item 재료원가, 노무원가 및 기타 제조원가 중 비정상적으로 낭비된 부분은 재고자산의 취득원가에 포함할 수 없다.
\end{choices}
\end{question}
\answer{2}
